\section{Architecture and Methodology}
\label{sec:method}

SecretGuard is architected as an autonomous multi-stage pipeline designed to execute seamlessly within developer Git pre-commit hooks and continuous integration workflows. Figure~\ref{fig:architecture} illustrates the five functional phases of the system.

\begin{figure*}[t]
\centering
\begin{tikzpicture}[
    node distance=1.2cm and 1.5cm,
    every node/.style={font=\small},
    block/.style={rectangle, draw=blue!70!black, fill=blue!5, thick, rounded corners=3pt, minimum width=3.8cm, minimum height=1.0cm, align=center},
    decision/.style={diamond, draw=orange!80!black, fill=orange!10, thick, aspect=1.8, inner sep=2pt, align=center},
    action/.style={rectangle, draw=green!60!black, fill=green!5, thick, rounded corners=3pt, minimum width=3.2cm, minimum height=0.9cm, align=center},
    route/.style={rectangle, draw=purple!70!black, fill=purple!5, thick, rounded corners=3pt, minimum width=3.2cm, minimum height=0.9cm, align=center},
    line/.style={draw, -latex, thick}
]

% Nodes
\node [block] (input) {\textbf{Staged Git Diff Stream}\\(\texttt{git diff -{}-cached})};

\node [block, below of=input] (ingest) {\textbf{Diff Streaming Ingestion}\\\textit{Filter additions (+), Ignore non-code}};

\node [block, below of=ingest] (det) {\textbf{Deterministic Extraction Engine}\\\textit{15+ Regex Rules, Shannon Entropy $H(X)$,}\\\textit{Pre-Filters \& Bounded Context $\le 200$ chars}};

\node [decision, below of=det, node distance=1.8cm] (scorer) {\textbf{Ambiguity Scorer}\\$S_{amb} = w_1 U_{reg} + w_2 \hat{H} + w_3 A_{ctx}$};

% Left branch: Unambiguous
\node [action, below left of=scorer, node distance=2.8cm, xshift=-0.5cm] (local) {\textbf{Local Resolution}\\\textit{Unambiguous Tier ($S < 0.2$)}\\\textbf{0ms Network Latency}};

% Right branch: Tiered LLM routing
\node [route, below right of=scorer, node distance=2.8cm, xshift=0.5cm] (router) {\textbf{Adaptive Model Gate}\\\textit{Profile A: Llama-3.1 8B ($0.2 \le S < 0.6$)}\\\textit{Profile B: GPT-4o-mini ($S \ge 0.6$)}};

\node [block, below of=scorer, node distance=3.6cm] (remediation) {\textbf{Autonomous Remediation Engine}\\\textit{1-Click Quarantine: Code Sanitization (\texttt{os.getenv} / \texttt{process.env}),}\\\textit{\texttt{.env} Migration, \texttt{.gitignore} Protection, SARIF v2.1.0 Export}};

% Connectors
\path [line] (input) -- (ingest);
\path [line] (ingest) -- (det);
\path [line] (det) -- (scorer);
\path [line] (scorer) -| node [near start, above] {$S_{amb} < 0.20$} (local);
\path [line] (scorer) -| node [near start, above] {$S_{amb} \ge 0.20$} (router);
\path [line] (local) |- (remediation);
\path [line] (router) |- (remediation);

\end{tikzpicture}

\caption{SecretGuard End-to-End Architectural Pipeline. Staged Git diff additions are streamed into the deterministic extraction engine. Candidates are evaluated via mathematical ambiguity scoring ($S_{amb}$), routing unambiguous tokens to instant local resolution and ambiguous candidates through privacy-preserved tiered model profiles.}
\label{fig:architecture}
\end{figure*}

\subsection{Diff Streaming Ingestion}
Traditional repository scanners load complete source files or walk entire commit trees, incurring severe I/O penalties. In contrast, SecretGuard's \texttt{GitDiffParser} streams the unified diff output generated by \texttt{git diff -{}-cached}. The parser operates as a memory-efficient generator:
\begin{itemize}
    \item It filters for file additions, extracting target paths from \texttt{diff -{}-git} and \texttt{+++ b/...} headers.
    \item It discards deleted lines (\texttt{-}) and unchanged surrounding context lines (\texttt{\phantom{+}}).
    \item Only added lines (\texttt{+}) are yielded as lightweight \texttt{DiffLine} structures containing the file path, line number, and raw string payload.
\end{itemize}
Concurrently, an ignore matcher evaluates paths against \texttt{.gitignore} and custom \texttt{.secretguardignore} rules, instantly pruning binary artifacts, minified bundles (\texttt{*.min.js}), lockfiles, and package directories.

\subsection{Deterministic Extraction and Context Bounding}
The deterministic extraction engine applies a multi-stage sieve to each incoming line:
\begin{enumerate}
    \item \textbf{Vendor Signature Library:} A collection of over 15 pre-compiled regular expressions subclassing \texttt{BaseRegexRule} targets structured credential formats, including AWS access and secret keys, GitHub classic and fine-grained personal access tokens, OpenAI project and legacy keys, Google API keys, Slack tokens, JWT structures, private key headers, Stripe keys, and database connection strings.
    \item \textbf{Candidate Pre-Filtering:} To suppress obvious false alarms before mathematical scoring, candidates are evaluated by \texttt{CandidateFilter}. This component suppresses canonical UUIDs, pure hex commit hashes, repeated uniform character sequences, and template placeholders (e.g., \texttt{YOUR\_API\_KEY}, \texttt{<TOKEN>}, \texttt{\$\{ENV\_VAR\}}).
    \item \textbf{Mathematical Shannon Entropy:} For unstructured or heuristic candidates, the engine calculates the Shannon entropy $H(X)$~\cite{shannon1948mathematical} over the character distribution:
    \begin{equation}
        H(X) = -\sum_{x \in \mathcal{A}} P(x) \log_2 P(x)
    \end{equation}
    where $\mathcal{A}$ denotes the unique alphabet in candidate $X$ of length $N$, and $P(x) = \frac{\text{count}(x)}{N}$. Strings with length $\ge 8$ and $H(X) \ge 4.2$ are flagged as high-entropy candidates.
    \item \textbf{Strict Context Window Bounding:} To eliminate data leakage to external verification services, surrounding code context is strictly bounded to $\le 200$ characters centered around the matched literal.
\end{enumerate}

\subsection{Mathematical Ambiguity Scoring Formulation}
A central innovation in SecretGuard is the closed-form mathematical ambiguity metric $S_{amb}(c) \in [0.0, 1.0]$ for candidate $c$:
\begin{equation}
    S_{amb}(c) = \text{clamp}\Big( w_1 \cdot U_{regex}(c) + w_2 \cdot \hat{H}(c) + w_3 \cdot A_{ctx}(c),\; 0.0,\; 1.0 \Big)
    \label{eq:ambiguity}
\end{equation}
where default calibrated weights are $w_1 = 0.4$, $w_2 = 0.3$, and $w_3 = 0.3$.

\subsubsection{Regex Uncertainty ($U_{regex}$)}
Quantifies the structural certainty of the matching rule:
\begin{equation}
    U_{regex}(c) = \begin{cases}
    0.0, & \text{if } c.\text{rule\_id} \in \text{STRUCTURED\_RULE\_IDS} \\
    1.0, & \text{otherwise (generic entropy or heuristic rule)}
    \end{cases}
\end{equation}

\subsubsection{Normalized Entropy ($\hat{H}$)}
Normalizes raw Shannon entropy across the active uncertainty range $[H_{min}, H_{max}] = [3.5, 6.0]$:
\begin{equation}
    \hat{H}(c) = \text{clamp}\left( \frac{H(c) - 3.5}{6.0 - 3.5},\; 0.0,\; 1.0 \right)
\end{equation}

\subsubsection{Contextual Ambiguity ($A_{ctx}$)}
Captures syntactic and semantic markers in the file path and surrounding code:
\begin{equation}
    A_{ctx}(c) = \begin{cases}
    1.0, & \text{if path or context contains mock/test indicators} \\
    0.2, & \text{if context matches explicit sensitive keywords} \\
    0.8, & \text{otherwise (generic identifiers or raw literals)}
    \end{cases}
\end{equation}
where mock indicators include \texttt{test}, \texttt{mock}, \texttt{fixture}, \texttt{dummy}, and \texttt{sample}, and sensitive keywords include \texttt{password}, \texttt{secret}, \texttt{api\_key}, and \texttt{access\_token}.

\subsection{Adaptive Routing Policy}
Based on $S_{amb}(c)$, the candidate is routed through one of three tiers:
\begin{itemize}
    \item \textbf{Unambiguous Tier ($S_{amb} < 0.20$):} Represents structured vendor credentials located in production code. The candidate resolves locally with zero network calls and 0\,ms network latency.
    \item \textbf{Low-Complexity Ambiguous Tier ($0.20 \le S_{amb} < 0.60$):} Represents partial ambiguity (e.g., a structured key inside a test file). Routed to \textbf{Model Profile A} (lightweight, ultra-fast model: Groq \texttt{llama-3.1-8b-instant}, max 150 tokens, 2.0s timeout).
    \item \textbf{High-Complexity Ambiguous Tier ($S_{amb} \ge 0.60$):} Represents high ambiguity (e.g., generic high-entropy strings embedded in complex test fixtures). Routed to \textbf{Model Profile B} (high-reasoning model: OpenRouter \texttt{openai/gpt-4o-mini}, max 200 tokens, 2.0s timeout).
\end{itemize}

\subsection{Autonomous Remediation and Quarantine Engine}
When a credential is confirmed, SecretGuard activates \texttt{SecretQuarantineManager} to perform atomic remediation:
\begin{enumerate}
    \item \textbf{Source Code Sanitization:} Analyzes target file language by extension. Replaces raw string literals in Python with \texttt{os.getenv("<VAR\_NAME>")} and in JavaScript/TypeScript with \texttt{process.env.<VAR\_NAME>}.
    \item \textbf{Environment Variable Migration:} Infers standard snake-case variable names from rule identifiers (e.g., \texttt{SG-AWS-001} $\to$ \texttt{AWS\_ACCESS\_KEY\_ID}) and appends the variable assignment to \texttt{.env}.
    \item \textbf{Git Exclusion Safeguard:} Checks repository \texttt{.gitignore}. If \texttt{.env} is not ignored, appends it automatically to ensure credentials are never committed.
\end{enumerate}
